\label{sec:systemdesc}

This section provides a consolidated \textbf{subsystem definition} and \textbf{data-flow} narrative for the AgriHome Vision Console, including the Summer~2026 Raspberry~Pi / Klipper edge path. Figure~\ref{fig:subsys-flow} maps major building blocks to the five layers introduced in Section~2. The text following the figure walks through each layer and explains how information moves from user or agent intent to durable storage and back.

\begin{figure}[H]
\centering
\scalebox{0.86}{
\begin{tikzpicture}[
  font=\footnotesize,
  arr/.style={-{Stealth}, thick},
  band/.style={draw, rounded corners, align=left, inner sep=7pt, text width=13.0cm},
]
\node[band, fill=leafgreen!16] (r1) {
  \textbf{Presentation Layer}\par
  Auth UI \textbullet{} Dashboard \textbullet{} Tray console \textbullet{} Devices UI \textbullet{} Tray Raspberry~Pi panel / Take Picture \textbullet{} Plant detail \textbullet{} Mesh \textbullet{} Schedule editor \textbullet{} PWA shell
};
\node[band, fill=blue!7, below=0.45cm of r1] (r2) {
  \textbf{Application Layer}\par
  HTTP/GraphQL entry \textbullet{} Session \& device-key auth \textbullet{} \texttt{/api/raspberry-pi/*} \textbullet{} \texttt{/api/devices/*} \textbullet{} Domain orchestration \textbullet{} optional streamer still client \textbullet{} Response assembly
};
\node[band, fill=orange!12, below=0.45cm of r2] (r3) {
  \textbf{Processing Layer}\par
  Species/leaf inference \textbullet{} Tray vision adapter \textbullet{} Edge disease hook \textbullet{} Capture schedule runner \textbullet{} Result normalization
};
\node[band, fill=gray!13, below=0.45cm of r3] (r4) {
  \textbf{Data Layer}\par
  PostgreSQL (\texttt{edge\_devices}, commands, poses, schedules, \texttt{camera\_captures}, \texttt{plants}, \texttt{prediction\_results}, \ldots) \textbullet{} Image/media object store
};
\node[band, fill=teal!14, below=0.45cm of r4] (r5) {
  \textbf{Hardware Layer}\par
  Raspberry~Pi \textbullet{} Agri-Home/klipper + \texttt{fswebcam} (\texttt{save\_image.sh}) \textbullet{} optional HTTP streamer \textbullet{} \texttt{agrihome\_agent}
};
\draw[arr] (r1.south) -- (r2.north);
\draw[arr, dashed] (r2.south) -- node[right, font=\scriptsize] {invoke when required} (r3.north);
\draw[arr] (r3.south) -- (r4.north);
\draw[arr, bend left=12] (r2.west) to node[left, font=\scriptsize, sloped, pos=0.4] {CRUD} (r4.west);
\draw[arr, bend left=16] (r5.west) to node[left, font=\scriptsize, align=center] {register /\\heartbeat / ingest} (r2.west);
\draw[arr, dashed, bend right=10] (r2.east) to node[right, font=\scriptsize] {snapshot} (r5.east);
\end{tikzpicture}
}
\caption{Subsystem definitions and primary data flow (edge hardware talks to the application layer; application also reads/writes the data layer directly for most CRUD paths).}
\label{fig:subsys-flow}
\end{figure}

\subsection{Presentation layer composition}

The Presentation Layer is organized around \textbf{workspaces} that match operator mental models. \textbf{Authentication and registration} establish identity and tie the browser to a server session. The \textbf{dashboard} summarizes tray health and recent activity. The \textbf{tray console} shows imagery, monitoring charts, and plant rows for a single tray, including a \textbf{Raspberry~Pi panel} with Take Picture, device link, and Klipper/streamer URL edit. The \textbf{Devices} view lists provisioned edge devices (serial, status, last heartbeat, linked tray). \textbf{Plant detail} focuses on one organism. \textbf{Mesh views} aggregate trays. The \textbf{schedule editor} binds capture policy to trays or meshes, including \texttt{raspberry-pi-edge} destinations. The \textbf{PWA shell} supplies offline-friendly installation where supported.

User gestures in these areas become HTTPS requests to the Application Layer; no presentation component queries the database directly.

\subsection{Application layer composition}

The Application Layer terminates TLS, verifies cookies/tokens or \textbf{device API keys}, and maps the caller to an account or device context. \textbf{Validation} rejects malformed payloads before any side effect. \textbf{Domain orchestration} implements use cases: plant/tray CRUD, meshes, schedules, monitoring, edge register/heartbeat/ingest, operator device actions (\texttt{capture}, \texttt{linkTray}, \texttt{revoke}, \texttt{rotateKey}, \texttt{updateKlipperUrl}), and optional HTTP streamer still attempts. When a use case requires inference, the same layer forwards work to the Processing Layer and merges outcomes into domain writes. \textbf{Response assembly} maps internal results to JSON or GraphQL shapes consumed by the UI or agent.

\subsection{Processing layer composition}

The Processing Layer exposes inference contracts and asynchronous capture helpers. A \textbf{species or leaf adapter} may classify a single-plant photo. A \textbf{tray vision adapter} may estimate plant count or regions. The \textbf{edge disease hook} optionally classifies captures after Pi ingest. The \textbf{capture schedule runner} evaluates due schedules with destination \texttt{raspberry-pi-edge} and enqueues \texttt{capture\_now} (optionally with \texttt{runPoses}). \textbf{Result normalization} keeps bounding boxes, labels, and scores consistent before persistence.

\subsection{Data layer composition}

The Data Layer splits \textbf{transactional} concerns from \textbf{bulk media}. Write coordination sequences inserts so relational rows and file metadata stay aligned. Integrity constraints guard keys and ownership. Query execution runs against \textbf{PostgreSQL}, which holds trays, plants, meshes, schedules, monitoring events, predictions, reports, and edge tables (\texttt{edge\_devices}, \texttt{edge\_device\_commands}, \texttt{capture\_pose\_sequences}, \texttt{capture\_poses}, \texttt{capture\_schedules}, \texttt{camera\_captures}). The \textbf{object store} retains images; the relational model stores pointers.

\subsection{Hardware layer composition}

The Hardware Layer is the bench capture plane: a \textbf{Raspberry~Pi} host, \textbf{Agri-Home/klipper} with \texttt{fswebcam}, optional HTTP streamer, and \textbf{\texttt{agrihome\_agent}}. Agents register and heartbeat; operators trigger Take Picture; schedules drive multi-pose walks. Primary capture is agent + \texttt{fswebcam}; optional server-side stills use a LAN-reachable \texttt{klipper\_url}.

\subsection{End-to-end data flow summary}

The data flow begins when a user interacts with the Presentation Layer or when an edge agent heartbeats. Requests pass to the Application Layer, which authorizes and validates them. Ordinary reads query the Data Layer. Take Picture primarily queues a command claimed by the Hardware Layer agent (\texttt{fswebcam}); it may optionally pull JPEG bytes from an HTTP streamer. Vision-assisted actions call the Processing Layer, then commit structured outputs through the Data Layer. Responses propagate back to the Presentation Layer (or agent) for rendering or further capture steps.
